\section{Two bounded operators on signed moment densities}
\label{ado:sec:operators}
For $f\in L^1(0,1)$ define its moments by
\[
 m_n(f)=\int_0^1u^{n-1}f(u)\dd u,\qquad n\ge1.
\]
We use the two operators
\begin{align}
 (\Aop f)(u)&=\int_u^1\frac{f(v)}v\dd v,
 \label{ado:eq:A}\\
 (\Bop f)(u)&=\int_0^u\frac{f(v)}{1-v}\dd v
              -\int_u^1\frac{f(v)}v\dd v.
 \label{ado:eq:B}
\end{align}
For every fixed $0<u<1$, both integrals in \eqref{ado:eq:B} are
absolutely convergent. In particular, this definition never subtracts
two divergent endpoint integrals. The operator $\Aop$ is the
outer-index-raising operator already present at depth two in the
manuscript~\cite{repoSigned}. The companion $\Bop$ encodes a strict
prefix sum of moments.

\begin{lemma}[Norm, derivative, and moment identities]
\label{ado:lem:operators}
Both operators map $L^1(0,1)$ to itself and satisfy
\begin{equation}\label{ado:eq:opnorm}
 \norm{\Aop f}_1\le\norm f_1,\qquad
 \norm{\Bop f}_1\le2\norm f_1.
\end{equation}
Their moments are
\begin{equation}\label{ado:eq:opmoments}
 m_n(\Aop f)=\frac{m_n(f)}n,\qquad
 m_n(\Bop f)=\frac1n\sum_{j=1}^{n-1}m_j(f).
\end{equation}
If $f$ is continuous in the open interval, then
\begin{equation}\label{ado:eq:opderivatives}
 (\Aop f)'(u)=-\frac{f(u)}u,\qquad
 (\Bop f)'(u)=\frac{f(u)}{u(1-u)}.
\end{equation}
Moreover $\Aop f$ extends continuously to one with value zero, and
$m_1(\Bop f)=0$.
\end{lemma}
\begin{proof}
Tonelli's theorem applied to absolute values gives
\[
 \int_0^1\int_u^1\frac{|f(v)|}{v}\dd v\dd u
 =\int_0^1|f(v)|\dd v.
\]
Similarly,
\[
 \int_0^1\int_0^u\frac{|f(v)|}{1-v}\dd v\dd u
 =\int_0^1|f(v)|\dd v.
\]
These identities prove the two operator bounds and justify every
following application of Fubini to a signed integrand. For the first
operator,
\[
 m_n(\Aop f)=\int_0^1\frac{f(v)}v
       \left(\int_0^v u^{n-1}\dd u\right)\dd v
 =\frac{m_n(f)}n.
\]
The moment of the first term of $\Bop f$ is
\[
 \frac1n\int_0^1\frac{1-v^n}{1-v}f(v)\dd v
 =\frac1n\sum_{j=1}^n m_j(f).
\]
Subtracting $m_n(f)/n$ yields the strict prefix in
\eqref{ado:eq:opmoments}. The derivative formulas follow locally from
the fundamental theorem of calculus. Finally,
$\int_u^1f(v)/v\dd v\to0$ as $u\uparrow1$, because $1/v$ is bounded
near one and $f$ is integrable there.
\end{proof}

\subsection{Construction at arbitrary depth}
The positive starting density is
\begin{equation}\label{ado:eq:seed}
 f_b(u)=\frac{(-\log u)^{b-1}}{\Gamma(b)},\qquad b>0.
\end{equation}
Its mass is one and the Gamma integral gives $m_n(f_b)=n^{-b}$.
Define, with composition read from right to left,
\begin{equation}\label{ado:eq:kappaconstruction}
 \kappa_{\svec}
 =\Aop^{s_1-1}\Bop\Aop^{s_2-1}\Bop\cdots
       \Aop^{s_{d-1}-1}\Bop f_b.
\end{equation}
There are $d-1$ occurrences of $\Bop$; all powers of $\Aop$ are
nonnegative integer powers. At depth one the right side is simply $f_b$.

\begin{proposition}[Moment realization and continuation]
\label{ado:prop:realization}
The density \eqref{ado:eq:kappaconstruction} is smooth on $(0,1)$,
belongs to $L^1(0,1)$, and satisfies
\begin{equation}\label{ado:eq:moments}
 m_n(\kappa_{\svec})=c_{\svec}(n),\qquad
 \norm{\kappa_{\svec}}_1\le2^{d-1}.
\end{equation}
The principal continuation is
\begin{equation}\label{ado:eq:signedtransform}
 L_{\svec}(z)=z\int_0^1\frac{\kappa_{\svec}(u)}{1-zu}\dd u,
 \qquad z\in\Om.
\end{equation}
The density is the unique finite signed measure realizing these moments.
\end{proposition}
\begin{proof}
The coefficient recurrence for a nonempty tail $\boldsymbol{t}$ is
\[
 c_{(a,\boldsymbol{t})}(n)
 =\frac1{n^a}\sum_{j=1}^{n-1}c_{\boldsymbol{t}}(j).
\]
It is exactly the effect of $\Aop^{a-1}\Bop$ in
\eqref{ado:eq:opmoments}. Induction proves the moments and the norm
bound. The derivative identities prove local smoothness, starting from
\eqref{ado:eq:seed}. Expanding the denominator in the open disk is
justified by the integrable bound $|\kappa|/(1-|z|)$, and gives the
series \eqref{ado:eq:series}.

On a compact subset of $\Om$, $|1-zu|$ is bounded away from zero
uniformly in $u\in[0,1]$. The signed integral is consequently holomorphic
there and continues the same germ. If two finite signed measures have
these moments, their difference integrates every polynomial to zero.
Uniform polynomial approximation on $[0,1]$, followed by uniqueness of
a finite measure from its continuous-function integrals, proves equality.
\end{proof}

\begin{corollary}[Ordinary convergence on the nontrivial unit circle]
\label{ado:cor:boundaryseries}
For \eqref{ado:eq:domain}, $c_{\svec}(n)\to0$ and
\begin{equation}\label{ado:eq:coefficientvariation}
 \sum_{n\ge1}|c_{\svec}(n+1)-c_{\svec}(n)|
 \le\norm{\kappa_{\svec}}_1\le2^{d-1}.
\end{equation}
Thus the original series \eqref{ado:eq:series} converges at every
$|z|=1$, $z\ne1$, and agrees there with \eqref{ado:eq:signedtransform}.
\end{corollary}
\begin{proof}
Dominated convergence proves the coefficient limit. Also
\[
 |c(n+1)-c(n)|\le\int_0^1u^{n-1}(1-u)|\kappa(u)|\dd u.
\]
Summing the positive majorants gives \eqref{ado:eq:coefficientvariation}.
Summation by parts against the bounded partial sums of $z^n$ then proves
convergence when $|z|=1$, $z\ne1$. Abel's limiting theorem identifies its
sum with the radial limit of the disk series, which in turn equals the
continuous signed integral at that point.
\end{proof}
This conclusion concerns the parameter domain of this article. It is not
a claim that all of the manuscript's more general fractional-order
boundary series converge ordinarily.
